% Candidate statements and complete proofs for the DEG_L paper.
% Checked by /root/semireg_t7, GPT-6 Astra Pro; inherited effort not exposed.
% This is an input fragment, not a standalone article.
% Source locators are written explicitly so no guessed bibliography keys are used.

\subsection{The actual surface and its deformation problem}

\begin{proposition}[S-01: Allowed source covers do not preserve an embedding premise]
\label{prop:semireg-cover-domain}
Let $f:Y\to A$ be a surface morphism obtained from the four rational
Hodge maps in 032, Proposition~3.1. If $\pi:Y'\to Y$ is a connected finite
\'etale cover of degree $d>1$, the pulled-back maps still satisfy its
nonzero-period criterion, but $f\pi$ is not a closed immersion.
Moreover $(f\pi)_*[Y']=d f_*[Y]$.
\end{proposition}
\begin{proof}
Pullback preserves rational Hodge maps and multiplies the integral of
a top-degree class by $d$. Every geometric point of $Y$ has $d$ distinct
preimages with the same image under $f\pi$, so this morphism is not
injective on geometric points. The last assertion is the proper
pushforward formula for a finite map. The deeper principal covers
used here are explicitly permitted in 032, Corollary~6.7 and \S7.
\end{proof}

\begin{proposition}[S-02: The exact conormal gate]
\label{prop:semireg-dual-gate}
For a smooth closed immersion $i:Y\hookrightarrow A$ of a surface in
an abelian $g$-fold, $g\geq3$, write $N=N_{Y/A}$ and
$W=H^0(A,\Omega_A^1)$. Bloch semiregularity is equivalent to
surjectivity of
\[
 H^1(A,\mathcal O_A)\otimes\bigwedge^3W
 \longrightarrow H^1(Y,N^\vee\otimes K_Y).
\]
The map is restriction followed by the natural morphism
$\bigwedge^3 i^*\Omega_A^1\to N^\vee\otimes K_Y$.
\end{proposition}
\begin{proof}
Bloch's map is
$H^1(Y,N)\to H^{g-1}(A,\Omega_A^{g-3})$.
Serre duality identifies its dual source with $H^1(A,\Omega_A^3)$
and its dual target with $H^1(Y,N^\vee\otimes K_Y)$.
The conormal exact sequence induces the displayed exterior-power
map, since the quotient $\Omega_Y^1$ has rank two.
This is the defining dual map in Buchweitz--Flenner \S8.1;
Proposition~8.2 identifies it with Bloch's map.
Finally $\Omega_A^3=\mathcal O_A\otimes\bigwedge^3W$,
and finite-dimensional duality gives the equivalence.
\end{proof}

\begin{proposition}[S-03: A nonzero period does not force semiregularity]
\label{prop:semireg-period-counterexample}
There is a smooth surface in a CM abelian threefold with a nonzero
product of two holomorphic and two antiholomorphic one-forms for
which the Bloch semiregularity kernel has dimension at least eight.
\end{proposition}
\begin{proof}
Let $E$ be a CM elliptic curve, $B=E^2$, and
$L=\mathcal O_E(3[0])\boxtimes\mathcal O_E(3[0])$.
A general $C\in|L|$ is smooth and connected. Since $L^2=18$ and
$K_B=\mathcal O_B$, adjunction gives $g(C)=10$ and
$N_{C/B}=K_C$. Put $A=B\times E$ and $Y=C\times E$.
Then $N_{Y/A}=K_C\boxtimes\mathcal O_E$, so K\"unneth gives
$h^1(Y,N)=1+10=11$. The divisor semiregularity target
$H^2(A,\mathcal O_A)$ has dimension three.
The curve $C$ dominates the first elliptic factor: otherwise its
ample divisor class would be a fibre class. Hence a nonzero one-form
from that factor restricts to a nonzero form $\alpha$ on $C$.
For a nonzero one-form $\beta$ on the last $E$, the integral
$\int_Y(\alpha\wedge\beta)\wedge\overline{\alpha\wedge\beta}$
is nonzero, by the product formula and positivity of the two
one-dimensional integrals.
\end{proof}

\begin{proposition}[S-04: Rigid generating sources]
\label{prop:semireg-rigid-source}
Let $Y$ be a smooth projective surface with $H^1(Y,T_Y)=0$ and
$i:Y\hookrightarrow A$ a smooth closed immersion. Let $G\subset A$
be the abelian subvariety generated by a translate of $i(Y)$.
Then the ambient obstruction map satisfies
\[
\ker(\operatorname{ob}_i)=
\operatorname{Lie}(A)\otimes
\ker\bigl(H^1(A,\mathcal O_A)\to H^1(Y,\mathcal O_Y)\bigr).
\]
In particular it is injective when $G=A$. If in addition $i$ is
Bloch semiregular, this is also the infinitesimal Hodge-locus kernel
of the entire class $i_*[Y]$.
\end{proposition}
\begin{proof}
Restriction on $H^1(T_A)$ is $\operatorname{id}_{\operatorname{Lie}(A)}
\otimes i^*$, since $T_A$ is trivial. The normal sequence and
$H^1(T_Y)=0$ make $H^1(i^*T_A)\to H^1(N)$ injective.
This proves the kernel formula. The map $\operatorname{Alb}(Y)\to A$
has image $G$, so $i^*$ on $H^1(\mathcal O)$ is injective precisely
when $G=A$. Buchweitz--Flenner, Corollary~4.12 and Proposition~8.2,
identify
\[
 \sigma_i(\operatorname{ob}_i(\kappa))
 =\kappa\mathbin{\lrcorner}i_*[Y].
\]
If $\sigma_i$ is injective, the kernels of the two sides agree.
\end{proof}

The actual 032 source is a compact smooth complex two-ball quotient.
Calabi--Vesentini, \emph{Sur les vari\'et\'es complexes compactes
localement sym\'etriques}, Th\'eor\`eme~3, p.~315, gives
$H^1(Y,T_Y)=0$. The nonzero four-factor period does not establish
$G=A$. For example, separate injectivity on the two eigenspaces
of a quadratic action is insufficient: in
$A=E^n\times E^n$ with action $(a,\bar a)$, pullback to the diagonal
$E^n$ is injective on each eigenspace, but their images coincide.

\begin{proposition}[S-05: A quantitative obstruction for the actual morphism]
\label{prop:semireg-map-dimension}
Let $Y$ be a compact smooth complex two-ball quotient and
$f:Y\to A$ a morphism with two-dimensional image, where $\dim A=g$.
Let $\tau_f$ be the Buchweitz--Flenner morphism semiregularity map.
Then
\[
 \dim\ker\tau_f\ \geq\ gq(Y)-\binom{2g}{g-2}.
\]
\end{proposition}
\begin{proof}
In characteristic zero $df:T_Y\to f^*T_A=\mathcal O_Y^{\oplus g}$
is generically injective. Its kernel is a torsion-free sheaf of
generic rank zero, hence vanishes. With $N_f=\operatorname{coker}(df)$,
the derived dual of $L_{Y/A}$ is the two-term complex
$[T_Y\to f^*T_A]$ in degrees $0,1$, so
$T^2_{Y/A}(\mathcal O_Y)=H^1(Y,N_f)$.
Calabi--Vesentini and the cokernel sequence give
$H^1(\mathcal O_Y)^{\oplus g}\hookrightarrow H^1(N_f)$.
The target specified in Buchweitz--Flenner, Proposition~7.21,
Lemma~7.22 and the definition before Theorem~7.23, is
\[
 \bigoplus_{r=0}^{g-2}H^{r+2}(A,\Omega_A^r).
\]
Its dimension is
$\sum_{r=0}^{g-2}\binom gr\binom g{r+2}=\binom{2g}{g-2}$
by Vandermonde. Rank--nullity proves the assertion.
\end{proof}

\begin{corollary}[S-06: Failure at sufficiently deep allowed levels]
\label{cor:semireg-level-failure}
Fix a nonzero 032 surface construction $f:Y\to A$. There are
connected finite \'etale covers $\pi_j:Y_j\to Y$, preserving its
four pulled-back Hodge maps and nonzero period, for which
$\tau_{f\pi_j}$ is not injective. Such covers may be chosen among
a cofinal system of principal congruence levels.
\end{corollary}
\begin{proof}
The nonzero period implies $H^1(Y,\mathbf C)\ne0$.
Venkataramana, \emph{Virtual Betti numbers of compact locally
symmetric spaces}, Theorem~1, proves that finite-dimensional
direct-limit cohomology over finite covers equals the cohomology
of the compact dual. Here the compact dual is $\mathbf P^2$,
whose first cohomology is zero. Pullback under finite covers is
injective over $\mathbf Q$ by transfer, so the direct limit is
nonzero and therefore infinite-dimensional. Hence $q(Y_j)$ is
unbounded, and the preceding proposition applies with fixed $A$.
The congruence version is explicitly identified in the final Remark
on p.~3 of that paper. Its proof uses the same rational-conjugation
action, since rational conjugation and finite intersection preserve
congruence subgroups; principal subgroups are cofinal.
Persistence of the period follows from the first proposition.
\end{proof}

This is a statement about the deformation problem of a morphism.
The factorization
$\tau_f=\sigma_{Rf_*\mathcal O_Y}\circ\psi_f$
does not imply that a kernel of $\tau_f$ survives in the perfect
complex obstruction space: $\psi_f$ itself can have a kernel.
Nor does the corollary settle Bloch semiregularity of a particular
smooth reduced image.

\begin{proposition}[S-07: Descent of the cover obstruction space]
\label{prop:semireg-cover-invariants}
For a finite \'etale Galois cover $\pi:Y'\to Y$ with deck group $H$,
$N_{f\pi}=\pi^*N_f$ and
$H^1(Y',N_{f\pi})^H=H^1(Y,N_f)$.
\end{proposition}
\begin{proof}
Flat pullback identifies the tangent cokernel sequences.
Finite pushforward and the projection formula identify
$\pi_*\pi^*N_f=N_f\otimes\pi_*\mathcal O_{Y'}$.
In characteristic zero averaging over $H$ splits off precisely
the invariant summand $N_f$. Cohomology commutes with this direct
summand, giving the assertion.
\end{proof}

\subsection{Derived replacements and their remaining gate}

\begin{proposition}[S-09: An integral Chern-component projector]
\label{prop:semireg-k-projector}
Let $A$ be an abelian $g$-fold and $\xi\in K_0(A)$ with
$\operatorname{ch}_j(\xi)=0$ for $j<p$. For $m\geq2$ put
\[
 M=[m]^*,\qquad
 D=\prod_{j=p+1}^g(m^{2j}-m^{2p}),\qquad
 \eta=\prod_{j=p+1}^g(m^{2j}\operatorname{id}-M)\xi.
\]
Then $\eta\in K_0(A)$ and
$\operatorname{ch}(\eta)=D\operatorname{ch}_p(\xi)$ in rational
Betti cohomology.
\end{proposition}
\begin{proof}
The operator $M$ multiplies $H^{2j}(A,\mathbf Q)$ by $m^{2j}$.
Thus every higher Chern component is annihilated by its corresponding
factor, lower components were already zero, and the component of
index $p$ is multiplied by $D$. The polynomial has integer
coefficients. Integral K-classes are represented by perfect complexes
on the smooth projective variety, using finite locally free
resolutions and shifts to implement negative coefficients.
\end{proof}

Starting from an integral codimension-$p$ cycle of class $N_0t$,
one may use $\xi=\sum a_i[\mathcal O_{Z_i}]$; its leading Chern
character is the fundamental cycle. Alternatively, resolve each
$Z_i$ and use the proper pushforward of its structure sheaf;
Grothendieck--Riemann--Roch gives the same leading component.
Adding line-bundle classes gives the useful character
$Nt+\sum_a b_a e^{a\theta}$, which remains Hodge on the polarized
Route-2 family. This constructs a K-class, not a semiregular object.

\begin{proposition}[S-08: Multiplicity creates trace kernels]
\label{prop:semireg-multiplicity}
If $E$ is perfect and $\operatorname{Ext}^2(E,E)\ne0$, then
$E^{\oplus m}$ is not fully semiregular for any $m\geq2$.
\end{proposition}
\begin{proof}
Its degree-two self-Ext group is
$\operatorname{Ext}^2(E,E)\otimes\operatorname{End}(\mathbf C^m)$.
The Atiyah class is $\operatorname{at}(E)\otimes\operatorname{id}$;
therefore every semiregularity component is the corresponding
component for $E$ tensored with the matrix trace.
A nonzero degree-two Ext class tensored with a nonzero traceless
matrix lies in the kernel.
\end{proof}

\begin{proposition}[S-10: Full Fourier--Mukai transport]
\label{prop:semireg-full-fm}
For the fixed enhanced Fourier--Mukai equivalence in the supplied
032 companion \S5.2, the full trace maps satisfy
\[
 \sigma_{\Phi Q}\circ\Phi_{\rm Ext}
   =\Phi^H_{HH}\circ\sigma_Q
 \quad\hbox{on }\operatorname{Ext}^2_A(Q,Q).
\]
Thus full semiregularity is invariant under this equivalence.
\end{proposition}
\begin{proof}
Use the companion's enhanced adjunctions and notation.
Put
$b_Q(v)=\operatorname{Tr}_{Q,H_A[2]}(A_Q^H[2]\circ v)
\in HH_{-2}(A)$.
Equations (5.20), (5.24) identify its HKR image with $\sigma_Q(v)$.
For $\delta\in T^{-2}(A)$, trace-compatible adjunction gives
\[
 \langle b_Q(v),\delta\rangle
   =\operatorname{Tr}_{A,Q}(\delta_Q[2]\circ v).
\]
Indeed, if $\delta':H_A\to\omega_A[g-2]$ is the adjoint coefficient
map, move $\delta'[2]$ through the partial trace by (5.24);
evaluation then gives the right side. This is exactly (5.37)
with the displayed even shift.
The convolution proof of (5.33), with the same inverse adjunction
data and shifted by two, gives
$(\Psi^{-2}\delta)_{\Phi Q}=\theta_Q[-2]\Phi(\delta_Q)$.
By the trace preservation in (5.31),
\[
 \langle b_{\Phi Q}(\Phi_{\rm Ext}v),\Psi^{-2}\delta\rangle
  =\operatorname{Tr}_{A,Q}(\delta_Q[2]\circ v)
  =\langle b_Q(v),\delta\rangle.
\]
Definition (5.34) and nondegeneracy of the Serre pairing imply
$b_{\Phi Q}(\Phi_{\rm Ext}v)=\Phi_*^{HH}b_Q(v)$.
Apply (5.35) and the fixed HKR maps. The two vertical maps
are isomorphisms, proving the assertion about kernels.
This proves a refinement; Proposition~5.2 itself only asserted
transport on the obstruction-evaluation image.
\end{proof}

\begin{proposition}[S-11: An ample complete intersection need not be semiregular]
\label{prop:semireg-ci-counterexample}
Let $A$ be an abelian fourfold, $L$ very ample, and
$Y=D_1\cap D_2$ a smooth transverse intersection of two members
of $|L|$. Then $h^1(N_{Y/A})=24$, whereas
$h^3(A,\Omega_A^1)=16$. Thus its Bloch map has a kernel of
dimension at least eight.
\end{proposition}
\begin{proof}
The twisted Koszul sequence is
$0\to L^{-1}\to\mathcal O_A^{\oplus2}\to\mathcal I_Y\otimes L\to0$.
Kodaira vanishing and $K_A=\mathcal O_A$ give $H^j(A,L)=0$ for
$j>0$; Serre duality gives $H^j(A,L^{-1})=0$ for $j<4$.
Thus
$H^2(A,\mathcal I_Y\otimes L)=H^2(A,\mathcal O_A)^{\oplus2}$.
The restriction sequence gives
$H^1(Y,L|_Y)=H^2(A,\mathcal I_Y\otimes L)$, of dimension twelve.
Transversality identifies $N_{Y/A}=(L|_Y)^{\oplus2}$.
The Bloch target has dimension
$\binom43\binom41=16$, so rank--nullity finishes the proof.
\end{proof}

\begin{proposition}[S-12: A pure Weil character obstruction for coherent sheaves]
\label{prop:semireg-pure-weil-no-go}
Let $A_0$ be a projective member of Voisin's four-dimensional
marked $\mathbf Q(i)$-Weil torus family, and let $t\ne0$ be a
rational Weil class. No fully Buchweitz--Flenner-semiregular
coherent sheaf $F$ on $A_0$ satisfies
\[
 \operatorname{ch}_1(F)=\operatorname{ch}_3(F)=0,\qquad
 \operatorname{ch}_2(F)=t.
\]
The rank and top-degree character are unrestricted.
\end{proposition}
\begin{proof}
The unpolarized Weil torus parameter space is smooth near $A_0$;
Voisin, \emph{A counterexample to the Hodge conjecture extended
to K\"ahler varieties}, \S3, equations (3.1)--(3.8), describes it.
Lemma~3.1 shows that the rational Weil space remains of Hodge
type $(2,2)$. Hence the components with indices $I=\{1,2,3\}$
in the assertion remain horizontal Hodge classes. In dimension
four these are precisely the components tested by full
semiregularity.
Buchweitz--Flenner Theorem~5.1 and its proof on p.~179 would
therefore extend $F$ as a coherent flat sheaf on a neighborhood.
Voisin Proposition~3.2 and Proposition~2.1 show that very general
fibres in that neighborhood have $\operatorname{NS}=0$ and
$c_2(G)=0$ for every coherent analytic sheaf $G$.
On such a fibre $c_1(F)=c_2(F)=0$, whence
$\operatorname{ch}_2(F)=0$, contrary to the nonzero horizontal
class $t$.
\end{proof}

This last proposition is a coherent-sheaf statement for the
specified fourfold family. Extension to arbitrary perfect complexes
would require an analytic effectivity theorem for their formal
deformations; algebraic/formal Pridham alone is not that theorem.
Allowing a nonzero polarization term in the Chern character
avoids this particular contradiction.

\begin{proposition}[S-13: One exponential polarization term does not repair the coherent gate]
\label{prop:semireg-single-exponential-no-go}
In the fourfold setting of S-12, let $L$ be a polarization whose
Rosati involution restricts to conjugation on $\mathbf Q(i)$, and
write $\theta=c_1(L)$. If $a\in\mathbf Z$, $\lambda\in\mathbf Q^\times$,
and $t\ne0$ is a rational Weil class, no fully semiregular coherent
sheaf $F$ has
\[
 \operatorname{ch}(F)=r e^{a\theta}+\lambda t+c[\mathrm{pt}]
\]
for any $r,c\in\mathbf Q$.
\end{proposition}
\begin{proof}
Write $H^1(A_0,\mathbf C)=V_\sigma\oplus V_{\bar\sigma}$.
The Rosati condition places $\theta$ in
$V_\sigma\otimes V_{\bar\sigma}$, whereas the complexified Weil
space is
$\bigwedge^4V_\sigma\oplus\bigwedge^4V_{\bar\sigma}$.
Both eigenspaces have dimension four. Consequently
$\theta\smile t=0$, and hence $e^{-a\theta}t=t$.
Tensoring by $L^{-a}$ is a Fourier--Mukai autoequivalence.
The functoriality in Perry Remarks~2.3 and~2.5, p.~5, preserves
full semiregularity under this equivalence, as in S-10.
It preserves coherence as well.
For $F'=F\otimes L^{-a}$ the Chern-character formula gives
\[
 \operatorname{ch}(F')=r+\lambda t+c[\mathrm{pt}].
\]
Thus $\operatorname{ch}_1(F')=\operatorname{ch}_3(F')=0$ and
$\operatorname{ch}_2(F')=\lambda t\ne0$, contradicting S-12.
\end{proof}

The integer condition on $a$ ensures that this normalization is an
actual line-bundle twist. The conclusion does not exclude all
polarization polynomials: it excludes the displayed single
exponential, including an arbitrary additional top-degree term.

\begin{proposition}[S-14: An equivariant perfect-complex gate implies bounded division]
\label{prop:semireg-perry-bounded-division}
Let $\pi:A\to S$ be an abelian scheme over a smooth irreducible
complex algebraic variety of finite type, with a relatively very
ample line bundle $H$ and $\theta=c_1(H)$. Let
$t\in H^0(S,R^{2p}\pi_*\mathbf Q(p))$ be a flat section whose values
are Hodge classes, where $0\leq p\leq\dim(A/S)$.
Suppose that for one $s_0\in S(\mathbf C)$ there exist:
\begin{enumerate}
\item a finite subgroup $G\subset A_{s_0}\times A_{s_0}^{\vee}$
with a specified coherent action on $\operatorname{Perf}(A_{s_0})$
by translations and tensoring by line bundles;
\item a weakly $G$-semiregular perfect complex $E_0$ with
$\operatorname{Ext}^{i}(E_0,E_0)=0$ for all $i<0$;
\item an algebraic class $B_0\in H^2(A_{s_0},\mathbf Q(1))$,
a nonzero $\lambda\in\mathbf Q$, and $P\in\mathbf Q[x]$ such that
\[
 e^{B_0}\operatorname{ch}(E_0)=\lambda t_{s_0}+P(\theta_{s_0})
\]
in the full even rational cohomology.
\end{enumerate}
Then there is an integer $C\geq1$ such that, for every
$s\in S(\mathbf C)$, there are an integer $1\leq N_s\leq C$
and effective integral codimension-$p$ cycles $Z_s^+,Z_s^-$,
each of $H$-degree at most $C$, with
\[
 N_s t_s=\operatorname{cl}(Z_s^+)-\operatorname{cl}(Z_s^-).
\]
In particular, there are single integers $N,D\geq1$ for the entire
family such that $Nt_s$ admits such representatives of degrees
at most $D$ on every fibre.
A sufficient form of the second hypothesis is a specified
$G$-equivariant structure on $E_0$ for which the ordinary full
semiregularity map is injective on
$\operatorname{Ext}^2(E_0,E_0)^G$, together with the same vanishing
of negative self-Ext groups.
\end{proposition}
\begin{proof}
The class
$w=\lambda t+P(\theta)$ is a global flat section of the full even
cohomology local system, and every component has the required
Hodge type. It extends
$w_0=e^{B_0}\operatorname{ch}(E_0)$.
Apply Alexander Perry, \emph{The semiregularity theorem for
equivariant noncommutative varieties},
arXiv:2604.00511v2 (24 June 2026), Theorem~1.2, p.~3.
Its second conclusion states that $w_s$ is algebraic for every
$s\in S(\mathbf C)$. Since $P(\theta_s)$ is algebraic and
$\lambda\ne0$, each $t_s$ is algebraic.
DIV1(2), applied to the fixed smooth projective family and $H$,
now supplies a single bound $C$ for the signed integral complexity.
This is exactly the displayed conclusion.
For a common multiplier take $N=\operatorname{lcm}(1,\ldots,C)$
and multiply the cycles on fibre $s$ by $N/N_s$.
Their degrees are at most $D=NC$.
Finally, Perry Definition~2.6 and Lemma~2.7, pp.~5--6, show that
the stated invariant-injectivity criterion implies weak
$G$-semiregularity over $\mathbf C$.
\end{proof}

The gate requires the entire displayed Chern character to remain
Hodge, a coherent categorical action, and vanishing of negative
self-Ext groups. A virtual K-class alone supplies none of these
deformation properties. The theorem supplies algebraicity on
every fibre; no properness assertion for a moduli space of
perfect complexes is used. The bound obtained from DIV1 is
existential. Here $G$ is a subgroup of the target's translation
and degree-zero line-bundle group, not an arbitrary deck group
of the source surface.

\begin{proposition}[S-15: Canonical source-cover equivariance retains a full kernel]
\label{prop:semireg-cover-equivariance}
Let $f:Y\to A$ be a proper morphism from a smooth projective
variety to an abelian variety, and put $E=Rf_*\mathcal O_Y$.
\begin{enumerate}
\item If $\rho:Y'\to Y$ is a finite \'etale Galois cover with
group $K$, the canonical $K$-equivariant lift of
$E'=R(f\rho)_*\mathcal O_{Y'}$ for the trivial action of $K$
on $\operatorname{Perf}(A)$ can be semiregular only if $E$
is fully semiregular.
\item Suppose $\rho$ is obtained by base change from an isogeny
$q:A'\to A$. Set
$G=\ker(q^*:A^\vee\to(A')^\vee)$, acting on
$\operatorname{Perf}(A)$ by the coherently chosen tensor
line bundles in $q_*\mathcal O_{A'}$.
The canonical $G$-equivariant lift of
$q_*q^*E=R(f\rho)_*\mathcal O_{Y'}$ can be semiregular only if
$E$ is fully semiregular.
\end{enumerate}
These assertions concern the canonical lifts supplied by the
covers; they do not exclude every possible new equivariant
structure on the underlying objects.
\end{proposition}
\begin{proof}
For (1), finite pushforward and characteristic-zero averaging give
\[
 (E')^K
 =Rf_*(\rho_*\mathcal O_{Y'})^K
 =Rf_*\mathcal O_Y=E.
\]
Because the categorical action is trivial,
$\operatorname{Perf}(A)^K$ is the finite product of copies of
$\operatorname{Perf}(A)$ indexed by irreducible representations
of $K$. The factor belonging to the trivial representation of
the canonical lift is $E$. Self-Ext groups, Hochschild homology,
and the semiregularity map decompose by these factors.
Semiregularity of the lift therefore implies that of $E$.

For (2), let $R=q_*\mathcal O_{A'}$. The character decomposition
of the finite abelian torsor gives
$R=\bigoplus_{P\in G}P$, with multiplication isomorphisms
$P\otimes Q\cong PQ$. These provide the indicated coherent
tensor action. A $G$-equivariant object is equivalently an
$R$-module object; since $q$ is finite \'etale this identifies
$\operatorname{Perf}(A)^G$ with $\operatorname{Perf}(A')$.
Under this equivalence the canonical lift of $q_*q^*E$
corresponds to $q^*E$.
Flat base change for the square defining $Y'$ identifies the
underlying object with the stated source-cover pushforward.

The pullback map
\[
 q^*:\operatorname{Ext}^2_A(E,E)
 \longrightarrow
 \operatorname{Ext}^2_{A'}(q^*E,q^*E)
\]
is injective: by adjunction its target is
$\operatorname{Ext}^2_A(E,E\otimes R)$, and the unit
$\mathcal O_A\to R$ has the retraction
$(\deg q)^{-1}\operatorname{Tr}_q$.
Semiregularity is natural for the enhanced functor $q^*$,
by Perry Remark~2.3, p.~5. Thus pullback carries
$\ker\sigma_E$ injectively into $\ker\sigma_{q^*E}$.
Semiregularity of the canonical lift would make the latter
kernel zero, proving the assertion.
\end{proof}

For an ordinary source deck transformation $k$ one has
$f\rho k=f\rho$. Consequently an induced translation of the
target satisfying this equivariance must be the identity.
An ordinary deck action therefore supplies internal
automorphisms of the pushed-forward object; it does not by
itself specify the nontrivial target action required in S-14.
Part (2) treats a special case where target tensor symmetries
really do arise, and shows why their canonical lift still
does not remove a pre-existing full semiregularity kernel.

There is also an independent primary-source confirmation of S-10:
Perry Remarks~2.3 and~2.5, p.~5, give functoriality of the
semiregularity morphism on the entire mapping complex and identify
its degree-two part with Buchweitz--Flenner under HKR.
An equivalence induces isomorphisms on both mapping complexes and
Hochschild homology, so the full kernel is invariant. This argument
does not depend on a raw versus Todd-modified HKR convention;
for abelian varieties the Todd class is one.

\begin{proposition}[S-16: A polarization exponential realizes a K-class by a vector bundle]
\label{prop:semireg-kclass-bundle-realization}
Let $X$ be a smooth projective variety, let $L$ be ample, and let
$\xi\in K_0(X)$. There exist integers $m,r\geq1$ and a vector
bundle $V$ such that
\[
 [V]=\xi+r[L^{-m}]\quad\hbox{in }K_0(X).
\]
In particular, if $\operatorname{ch}(\xi)=\lambda t$, then
$\operatorname{ch}(V)=r e^{-m c_1(L)}+\lambda t$.
In the setting of S-13, with $\lambda t\ne0$, every vector bundle
with this character fails full semiregularity.
\end{proposition}
\begin{proof}
On a smooth projective variety a K-class is a difference
$\xi=[U]-[W]$ of vector bundles: take finite locally free
resolutions and group their even and odd terms.
For sufficiently large $m$, the vector bundle $W\otimes L^m$
is generated by finitely many global sections.
The precise ampleness criterion is Stacks Project,
Proposition~28.27.13(7)--(8), tag 01Q3.
Choose a surjection
\[
 0\longrightarrow K\longrightarrow\mathcal O_X^{\oplus r}
   \longrightarrow W\otimes L^m\longrightarrow0 .
\]
One can add a redundant zero section to ensure $r\geq1$.
The kernel $K$ is locally free, since the surjection onto a
vector bundle splits locally. After tensoring by $L^{-m}$,
the K-group relation is
$[K\otimes L^{-m}]=r[L^{-m}]-[W]$.
Thus $V=U\oplus(K\otimes L^{-m})$ has the required class.
The Chern-character formula follows by additivity, and the
last assertion is S-13 with $a=-m$ and $c=0$.
\end{proof}

This gives an actual coherent vector-bundle representative
after one exponential addition. It supplies neither full nor
equivariant semiregularity; in the indicated fourfold case the
full version is ruled out by S-13.
